% GENERATED FILE. Edit canonical lessons, then run npm run generate:books.
% canonical-source-hash: fnv1a64:720040e4d3dcd66d
% canonical-lessons: BN-C140-dharalo, BN-C140-bhota, BN-C140-tatka, BN-C140-basi, BN-C140-kacha

\chapter{Sharp, Blunt, Fresh, Stale, Raw}
\label{ch:bn-sharp-blunt-fresh-stale-raw}
\hlchaptermodality{140}

\begin{chapteropening}
\textbf{By the end of this chapter:} \emph{I can say sharp, blunt, fresh, stale and raw.}

Complete the last lesson of chapter 140: I can say sharp, blunt, fresh, stale and raw.
\end{chapteropening}

\section[\texorpdfstring{\bn{ধারালো} (dhārālo)}{ধারালো (dhārālo)}]{\bn{ধারালো} (dhārālo) — sharp}

\label{lesson:BN-C140-dharalo}



\begin{quote}
Before the new one: say the Bengali for straight, then the Bengali for crooked.
\end{quote}

\subsection*{You'll want to know: \bn{ধারালো}}
\textbf{\bn{ধারালো}} — \emph{dhārālo} — \textquotedblleft{}sharp\textquotedblright{}.

\begin{grammarlens}[title={What you've built}]
One more word for this chapter.
\end{grammarlens}

\subsection*{Your turn}
\begin{itemize}
\raggedright
  \item \emph{Say it:} \emph{dhārālo}
  \item \emph{Say it:} \emph{dhārālo}, once more
  \item \emph{Say it:} say \emph{dhārālo}
  \item \emph{Recall:} say the Bengali for straight, then the Bengali for crooked, then say \emph{dhārālo} again
\end{itemize}

\begin{tcolorbox}[breakable,colback=teal!4,colframe=teal!35!black,title={Before you move on}]
What does \bn{ধারালো} mean? (\textquotedblleft{}Sharp\textquotedblright{}.) Say it once more.
\end{tcolorbox}
\section[\texorpdfstring{\bn{ভোঁতা} (bhõtā)}{ভোঁতা (bhõtā)}]{\bn{ভোঁতা} (bhõtā) — blunt}

\label{lesson:BN-C140-bhota}



\begin{quote}
Before the new one: say the Bengali for crooked, then the Bengali for sharp.
\end{quote}

\subsection*{You'll want to know: \bn{ভোঁতা}}
\textbf{\bn{ভোঁতা}} — \emph{bhõtā} — \textquotedblleft{}blunt\textquotedblright{}.

\begin{grammarlens}[title={What you've built}]
One more word for this chapter.
\end{grammarlens}

\subsection*{Your turn}
\begin{itemize}
\raggedright
  \item \emph{Say it:} \emph{bhõtā}
  \item \emph{Say it:} \emph{bhõtā}, once more
  \item \emph{Say it:} say \emph{bhõtā}
  \item \emph{Recall:} say the Bengali for crooked, then the Bengali for sharp, then say \emph{bhõtā} again
\end{itemize}

\begin{tcolorbox}[breakable,colback=teal!4,colframe=teal!35!black,title={Before you move on}]
What does \bn{ভোঁতা} mean? (\textquotedblleft{}Blunt\textquotedblright{}.) Say it once more.
\end{tcolorbox}
\section[\texorpdfstring{\bn{টাটকা} (ṭāṭkā)}{টাটকা (ṭāṭkā)}]{\bn{টাটকা} (ṭāṭkā) — fresh}

\label{lesson:BN-C140-tatka}



\begin{quote}
Before the new one: say the Bengali for sharp, then the Bengali for blunt.
\end{quote}

\subsection*{You'll want to know: \bn{টাটকা}}
\textbf{\bn{টাটকা}} — \emph{ṭāṭkā} — \textquotedblleft{}fresh\textquotedblright{}.

\begin{grammarlens}[title={What you've built}]
One more word for this chapter.
\end{grammarlens}

\subsection*{Your turn}
\begin{itemize}
\raggedright
  \item \emph{Say it:} \emph{ṭāṭkā}
  \item \emph{Say it:} \emph{ṭāṭkā}, once more
  \item \emph{Say it:} say \emph{ṭāṭkā}
  \item \emph{Recall:} say the Bengali for sharp, then the Bengali for blunt, then say \emph{ṭāṭkā} again
\end{itemize}

\begin{tcolorbox}[breakable,colback=teal!4,colframe=teal!35!black,title={Before you move on}]
What does \bn{টাটকা} mean? (\textquotedblleft{}Fresh\textquotedblright{}.) Say it once more.
\end{tcolorbox}
\section[\texorpdfstring{\bn{বাসি} (bāsi)}{বাসি (bāsi)}]{\bn{বাসি} (bāsi) — stale}

\label{lesson:BN-C140-basi}



\begin{quote}
Before the new one: say the Bengali for blunt, then the Bengali for fresh.
\end{quote}

\subsection*{You'll want to know: \bn{বাসি}}
\textbf{\bn{বাসি}} — \emph{bāsi} — \textquotedblleft{}stale\textquotedblright{}.

\begin{grammarlens}[title={What you've built}]
One more word for this chapter.
\end{grammarlens}

\subsection*{Your turn}
\begin{itemize}
\raggedright
  \item \emph{Say it:} \emph{bāsi}
  \item \emph{Say it:} \emph{bāsi}, once more
  \item \emph{Say it:} say \emph{bāsi}
  \item \emph{Recall:} say the Bengali for blunt, then the Bengali for fresh, then say \emph{bāsi} again
\end{itemize}

\begin{tcolorbox}[breakable,colback=teal!4,colframe=teal!35!black,title={Before you move on}]
What does \bn{বাসি} mean? (\textquotedblleft{}Stale\textquotedblright{}.) Say it once more.
\end{tcolorbox}
\section[\texorpdfstring{\bn{কাঁচা} (kā̃chā)}{কাঁচা (kā̃chā)}]{\bn{কাঁচা} (kā̃chā) — raw, unripe}

\label{lesson:BN-C140-kacha}



\begin{quote}
Before the new one: say the Bengali for fresh, then the Bengali for stale.
\end{quote}

\subsection*{You'll want to know: \bn{কাঁচা}}
\textbf{\bn{কাঁচা}} — \emph{kā̃chā} — \textquotedblleft{}raw, unripe\textquotedblright{}.

\begin{grammarlens}[title={What you've built}]
One more word for this chapter.
\end{grammarlens}

\subsection*{Your turn}
\begin{itemize}
\raggedright
  \item \emph{Say it:} \emph{kā̃chā}
  \item \emph{Say it:} \emph{kā̃chā}, once more
  \item \emph{Say it:} say \emph{kā̃chā}
  \item \emph{Recall:} say the Bengali for fresh, then the Bengali for stale, then say \emph{kā̃chā} again
\end{itemize}

\begin{tcolorbox}[breakable,colback=teal!4,colframe=teal!35!black,title={Before you move on}]
What does \bn{কাঁচা} mean? (\textquotedblleft{}Raw, unripe\textquotedblright{}.) Say it once more.
\end{tcolorbox}
